\begin{apartat}{1}
Un mètode per a determinar les masses d'ions pesants consisteix a mesurar el temps que necessiten per a fer un nombre determinat de voltes en un camp magnètic conegut. En un d'aquests mesuraments, un ió amb una càrrega igual a la d'un electró fa 7,00 voltes en $1,29\ \mathrm{ms}$ en un camp magnètic perpendicular a la velocitat i amb un mòdul de $45,0\ \mathrm{mT}$. Feu una representació de la trajectòria de l'ió i dibuixeu en dues posicions d'aquesta trajectòria el vector força que actua sobre l'ió. Calculeu la massa de l'ió.
\end{apartat}

\begin{apartat}{1}
Un protó que es mou a una velocitat de $5,00\times 10^{5}\ \mathrm{m\,s^{-1}}$ entra en una regió de l'espai on hi ha un camp magnètic. El mòdul de la força que produeix el camp magnètic sobre la càrrega és $8,00\times 10^{-14}\ \mathrm{N}$. Calculeu el mòdul del camp magnètic. Especifiqueu clarament la direcció i el sentit d'aquest camp magnètic si les direccions i els sentits, tant de la força com de la velocitat, són els representats en la figura.

\figura[0.22]{fig1.png}
\end{apartat}

\dades{%
  Càrrega de l'electró, $q_\mathrm{e} = -1,60\times 10^{-19}\ \mathrm{C}$.\\
  Càrrega del protó, $q_\mathrm{p} = 1,60\times 10^{-19}\ \mathrm{C}$.}
